%-------------------------------------------------------------------------------
%	ABSTRACT PAGE
%-------------------------------------------------------------------------------

\addchaptertocentry{Хураангуй} % Хураангуйг гарчигт нэмэх

\begin{center}
{\scshape\Large \univname\par} % Их сургуулийн нэр
{\scshape\large \facname\par}\vspace{0.5cm} % Их сургуулийн нэр
{\huge\textbf{{Хураангуй}} \par}
\bigskip
{\Large{\ttitle} \par} % Тезисийн нэр
\bigskip

{\normalsize \shortname \par} % Зохиогчийн нэр
\addressname
\end{center}

\textit{\textbf{Түлхүүр үгс: \keywordnames}}
\bigskip

Орчин үеийн хөдөлмөрийн зах зээлд нэг ажлын байранд олон зуун CV ирдэг бөгөөд уламжлалт гараар хийгддэг сонгон шалгаруулалтын процесс нь цаг хугацаа их шаарддаг, хүний субъектив алдаанд өртөмтгий байдаг. Энэхүү дипломын ажилд уг асуудлыг шийдвэрлэх зорилгоор ажил горилогчийн ур чадварыг ажлын байрны шаардлагад тулгуурлан хиймэл оюунаар тайлбарлан үнэлэх веб суурьтай систем боловсруулсан.

Системийн судалгааны хэсэгт Монголын зах зээлд үйл ажиллагаа явуулдаг zangia .mn болон lambda.global системүүдийг судалж, машин сургалт, байгалийн хэл боловсруулалт болон томоохон хэлний загвар (LLM)-д суурилсан алгоритмуудыг харьцуулан шинжилсний үндсэн дээр тайлбарлах боломжтой (explainable) үнэлгээ гаргах чадвартай LLM-д суурилсан аргыг сонгосон.

Хөгжүүлсэн систем нь ажил олгогч болон ажил горилогч гэсэн хоёр үндсэн хэрэглэгчид зориулагдсан. Ажил горилогч PDF хэлбэрийн CV байршуулахад Anthropic компанийн Claude Sonnet загвар ажлын байрны шаардлагатай харьцуулан нийцлийн оноо, тохирсон болон дутуу ур чадварууд, шаардлага болон үүрэг хариуцлага бүрийн биелэлтийн дүн, сайжруулах зөвлөмжүүдийг агуулсан бүтэцтэй JSON үнэлгээ автоматаар гаргадаг. Ажил олгогч тал ирсэн CV-үүдийг AI үнэлгээний дүнгээр эрэмбэлэн харж, ажил горилогчид ур чадварын практик даалгавар илгээх, гүйцэтгэлийг үнэлэх, Google Calendar-тай интеграцтайгаар ярилцлагын цаг товлох боломжтой. Чатботын модуль нь PostgreSQL өгөгдлийн сангаас бүх идэвхтэй ажлын байрны мэдээллийг Claude API-д контекст болгон шууд дамжуулдаг тул хэрэглэгчийн асуултад тухайн үеийн бодит мэдээлэлд суурилсан хариулт өгнө.

Системийг Go (Golang) хэлэн дээр Chi router болон GORM-ийг ашиглан хөгжүүлсэн бөгөөд хэрэглэгчийн интерфэйсийг Nuxt 4 болон shadcn/vue UI сангаар, өгөгдлийн санд PostgreSQL-ийг ашигласан. AI үнэлгээний үр дүнг WebSocket дамжуулан бодит цагт хэрэглэгчид хүргэдэг. Боловсруулсан систем нь сонгон шалгаруулалтын процессыг автоматжуулж, тайлбарлах боломжтой AI үнэлгээгээр дамжуулан ажил олгогч болон ажил горилогч хоёрт нэгдсэн платформ дотор цогц үйлчилгээ үзүүлдэг.




